\begin{claim}{D2.1}{Proved}{The entire first tangent cohomology}
The Leray class $\xi\otimes[\gamma]$ generates $H^1(X,T_X)$.
The Kodaira--Spencer map of $c\mapsto X_c$ is an isomorphism
$T_c\mathcal P\to H^1(X_c,T_{X_c})$ at every admissible $c$.
\end{claim}
\begin{proof}
All line bundles other than the trivial summands in
$R^q\bar f_*\mathcal V=E\otimes R^q\bar f_*\OO_X$ are acyclic:
they have weights $-1,-2$. Hence $H^1(X,\mathcal V)=
\C(\xi\otimes\gamma)$ and its other positive-degree groups vanish.
The three direct images of $\bar f^*\OO(-5)$ have weights
$-5$, $(-5,-6)$, and $-6$, all acyclic by D1.1. Consequently (D.3)
induces an isomorphism in every cohomology degree. Sequence (D.10),
whose quotient has no cohomology in degrees zero and one, then gives
the claimed isomorphism $H^1(\mathcal V)\to H^1(T_X)$.

We use the Cech convention in which the derivative of a transition
arrow is composed with the inverse of the central arrow. A lattice
transition is translation by $\Pi_c\lambda$, and its derivative in
$c$ is $\gamma(\lambda)\xi$. The linear deck matrices are independent
of $c$. On the finite charts the affine translations contribute
$\gamma(v_j)\xi/m_j$; their $m_j$th iterates give
$\gamma(v_j)\xi$, precisely the lattice cocycle. On the overlap with
the smooth part the derivative of the logarithmic translating section
is $\gamma(v_j)(\log s_j)\xi/(2\pi i)$. Its Cech coboundary accounts
for the difference of these representatives. On the cusp the deck
multiplier differentiates to $\gamma(\lambda)\xi$ in logarithmic
coordinates, and the exponential overlap identifies it with the same
cocycle. Thus the global Kodaira--Spencer class is
$[\xi\otimes\gamma]$, with a minus sign if the tangent-lift convention
$[\bar\partial h]$ is used instead. The sign is recorded; nonvanishing
does not depend on it.

Nonvanishing can also be checked on a regular torus. In the quotient
description of $H^1(\OO)$ used above, $[\gamma]\ne0$; the first
coefficient of (D.5) is zero, so $\xi\gamma$ cannot be a scalar
multiple of the base Kodaira--Spencer class. Its image in
$H^1(T_X|_F)$ therefore remains nonzero. Since $H^1(X,T_X)$ is
one-dimensional by D1.3, this proves the asserted isomorphism.
\end{proof}
